\documentclass[11pt]{article}
\usepackage[margin=1in]{geometry}
\usepackage{amsmath,amssymb,amsthm}
\usepackage{array}
\usepackage{parskip}
\usepackage{booktabs}

\newtheorem{theorem}{Theorem}
\newtheorem{lemma}[theorem]{Lemma}
\newtheorem{corollary}[theorem]{Corollary}
\newtheorem{proposition}[theorem]{Proposition}
\theoremstyle{definition}
\newtheorem{definition}[theorem]{Definition}
\theoremstyle{remark}
\newtheorem{remark}[theorem]{Remark}

\newcommand{\F}{\mathbb{F}}
\newcommand{\PG}{\mathrm{PG}}
\newcommand{\floor}[1]{\left\lfloor #1 \right\rfloor}

\title{Disproof of conjecture \texttt{00000001003}}
\author{earthking11}
\date{2026-09-14}

\begin{document}
\maketitle

\section*{Verdict: FALSE}

Conjecture \texttt{00000001003} is false. We disprove it at the very first odd
prime power, $q=3$: the conjecture predicts a maximal cap size of
\[
q^{2}+q+1+\floor{(q+1)/3}\Big|_{q=3}=9+3+1+\floor{4/3}=14,
\]
but $\PG(4,3)$ contains an explicit cap of $20$ points, and $20>14$. The failure
is not an artefact of $q=3$: the same formula is beaten in $\PG(4,5)$ (an
explicit $47$-point cap against the predicted $33$), the even-$q$ clause fails
at $q=2$ (a $16$-point cap against the predicted $10$), and the description
``attained by an elliptic quadric with an explicit three-point augmentation'' is
internally inconsistent with the stated formula. The refutation is
unconditional; no unproved hypothesis is used.

\section{The conjecture, quoted verbatim}

From \texttt{conjectures/00000001003.md}:

\begin{quote}
\textbf{English.} Definition: A cap is a point set of $\PG(n,q)$ with no three
collinear. Conjecture: The maximal cap size in $\PG(4,q)$ is
$q^{2}+q+1+\floor{(q+1)/3}$ for all odd $q$, attained by an elliptic quadric
with an explicit three-point augmentation; for even $q$ the value increases by
$2$.
\end{quote}

The file is bilingual; the Chinese original is reproduced verbatim in
\texttt{README.md} (this PDF does not embed CJK glyphs). The two versions are
identical in content.

\begin{definition}[cap]
A \emph{cap} in $\PG(n,q)$ is a set of points no three of which are collinear.
Its \emph{size} is its cardinality; it is \emph{inclusion-maximal} if no further
point can be added while remaining a cap.
\end{definition}

Throughout, ``maximal'' is read as \emph{maximum} (the largest size), which is
the only reading under which the conjecture states a numerical value. We also
test the inclusion-maximal reading in Section~\ref{sec:readings}.

\section{The explicit 20-point cap in $\PG(4,3)$}

We use homogeneous coordinates on $\PG(4,3)$; every point is written in
normalised form, i.e.\ its first nonzero coordinate is $1$. The following
$20$ points form a cap:
\begin{center}
\small
\begin{tabular}{lllll}
\toprule
$(1,0,0,0,2)$ & $(0,0,0,1,2)$ & $(1,2,2,1,1)$ & $(0,1,0,0,0)$ & $(1,0,2,1,2)$\\
$(1,0,1,2,0)$ & $(0,1,0,1,0)$ & $(1,2,1,2,1)$ & $(1,1,2,2,0)$ & $(1,2,0,1,2)$\\
$(1,2,0,1,0)$ & $(1,2,1,0,2)$ & $(1,0,1,1,0)$ & $(0,1,1,1,2)$ & $(0,0,1,0,0)$\\
$(0,1,2,2,0)$ & $(1,0,2,2,2)$ & $(1,2,1,2,2)$ & $(0,1,2,2,2)$ & $(1,1,0,1,1)$\\
\bottomrule
\end{tabular}
\end{center}

That is, no three of these $20$ points are collinear. Since the conjecture
predicts $14$ at $q=3$, this already gives
\[
M_{2}(5,3)\;\ge\;20\;>\;14\;=\;q^{2}+q+1+\floor{(q+1)/3}\big|_{q=3},
\]
so the $q=3$ instance of the quantified claim ``for all odd $q$'' is false.

\subsection{Three independent verifications}

All three were carried out and agree on the same answer, $0$ collinear
triples. The script \texttt{reproduce.py} performs the first two from scratch.

\begin{enumerate}
\item \textbf{Line membership.} $\PG(4,3)$ has $121$ points and $1210$ lines;
  each line has $q+1=4$ points. Enumerating all lines and intersecting them with
  the $\binom{20}{3}=1140$ triples gives $0$ collinear triples.
\item \textbf{Rank over $\F_{3}$.} Three projective points are collinear iff the
  rank of their three coordinate vectors is $\le 2$. Gaussian elimination over
  $\F_{3}$ on all $1140$ triples gives $0$ triples of rank $\le 2$.
\item \textbf{Kernel-checked Lean.} The executable checker in
  \texttt{lean4/Main.lean} (core Lean, no Mathlib, no \texttt{sorry}) returns
  \texttt{true} on the list; the theorem \texttt{pts20\_isCap} is proved by
  \texttt{decide} and depends on no axioms.
\end{enumerate}

\section{The exact value $20$ (SAT-certified upper bound)}

We have proved $M_{2}(5,3)\ge 20$. The matching upper bound $M_{2}(5,3)\le 20$
(i.e.\ no $21$-point cap exists) is a finite computation, which we report as
\emph{SAT-certified} rather than kernel-checked. The argument is:

\begin{proposition}[upper bound, SAT certificate]
$M_{2}(5,3)\le 20$.
\end{proposition}

\begin{proof}[Argument]
Let $C\subseteq\PG(4,3)$ be a cap with $|C|\ge 21$. Any set of $\ge 21$ points
in $\PG(4,3)$ spans $\PG(4,3)$, hence contains $5$ linearly independent points.
The group $\mathrm{GL}(5,3)$ acts transitively on ordered bases of $\F_{3}^{5}$,
so after a projectivity we may assume that the $5$ standard basis points (a
partial frame) lie in $C$. The remaining points are then searched subject to the
constraint that no three chosen points are collinear; with the frame fixed and
$|C|\ge 21$ required, the resulting CNF is UNSAT (checked with a SAT solver
using a cardinality encoding). Therefore no such $C$ exists.
\end{proof}

This is a computer-assisted certificate over a finite search space, not a
kernel-checked proof; it is recorded here for completeness. The disproof of the
conjecture does \emph{not} depend on it: the lower bound $M_{2}(5,3)\ge 20>14$
already refutes the value $14$. The value $20$ agrees with the classical table
entry $M_{2}(5,3)=20$.

\section{The general formula fails at $q=5$}

The conjecture quantifies over all odd $q$, so failing at $q=3$ already
disproves it. It is nevertheless false at the next odd prime power as well, and
by a wide margin. In $\PG(4,5)$ ($781$ points) an explicit cap of $47$ points
exists; the conjecture predicts
\[
25+5+1+\floor{6/3}=33 .
\]
The $47$ points are exhibited explicitly in \texttt{reproduce.py} and verified
there by two methods (pairwise line membership and rank over $\F_{5}$): all
$\binom{47}{3}=16215$ triples have rank $3$, so the cap is valid. Hence
$M_{2}(5,5)\ge 47>33$. (A cap of size $46$ at $q=5$ was found earlier; the
deterministic list shipped in \texttt{reproduce.py} has $47$ points and is
strictly stronger.)

\section{The even-$q$ clause fails at $q=2$}

For even $q$ the conjecture adds $2$ to the formula, giving at $q=2$
\[
4+2+1+\floor{3/3}+2 = 10 .
\]
But $\PG(4,2)$ ($31$ points) contains the $16$-point cap
\[
\{\,(1,v)\;:\;v\in\F_{2}^{4}\,\},
\]
and in fact $M_{2}(4,2)=16=2^{4}$, the classical value $M_{2}(n,2)=2^{n}$. The
lower bound $16$ is verified in \texttt{reproduce.py}; the upper bound
$M_{2}(4,2)\le 16$ is again a finite SAT certificate ($\ge 17$ is UNSAT),
reported as such. Thus the even-$q$ clause is off by $6$ at $q=2$: not merely
does the formula plus $2$ fail, it fails in the opposite direction (the true
value is larger, not smaller).

\section{Internal inconsistency of the description}

The conjecture says the value is ``attained by an elliptic quadric with an
explicit three-point augmentation''. An elliptic quadric in this setting
contributes $q^{2}+1$ points, and augmenting it by three points gives
\[
q^{2}+4 .
\]
For $q=3$ this is $13$, whereas the stated formula gives $14$; in general
\[
q^{2}+q+1+\floor{(q+1)/3} - (q^{2}+4)
   \;=\; q-3+\floor{(q+1)/3} \;>\; 0 \qquad (q\ge 3 \text{ odd}).
\]
So the construction does not attain the stated value (and the stated value is
not produced by the construction). The two halves of the sentence are
inconsistent with each other. The same mismatch occurs at $q=5$: $29$ from the
quadric description versus $33$ from the formula.

\section{Alternative readings}
\label{sec:readings}

We tried to rescue the conjecture by varying the reading; none succeeds.

\begin{center}
\begin{tabular}{p{0.36\textwidth} p{0.56\textwidth}}
\toprule
Reading & Outcome \\
\midrule
Relax ``cap'' (allow some collinear triples) &
  Contradicts the file's own definition of a cap and has no canonical value;
  the conjecture then states no definite number, so it is not a well-posed
  claim. \\
``Maximal'' means inclusion-maximal &
  The $20$-point cap above is inclusion-maximal (verified: no point of
  $\PG(4,3)$ can be added to it). A survey of $3000$ random greedy
  inclusion-maximal caps produced sizes $16$--$20$ (distribution
  $16\!:\!84$, $17\!:\!292$, $18\!:\!1889$, $19\!:\!459$, $20\!:\!276$),
  \emph{never} $14$. So this reading does not force the value $14$ either. \\
``The formula is meant for $q\ge 5$'' &
  Fails at $q=5$: an explicit $47$-point cap versus the predicted $33$
  (Section~5). \\
Even-$q$ clause taken as the base formula &
  Fails at $q=2$: $16$ versus $10$ (Section~6). \\
\bottomrule
\end{tabular}
\end{center}

The first line is the only one that changes the meaning of a term rather than a
quantifier range; under the file's own definition of ``cap'' it does not apply.

\section{Reproducibility}

\texttt{reproduce.py} (Python~3 standard library only, no third-party packages)
rebuilds $\PG(4,3)$, checks $121$ points and $1210$ lines, verifies the
$20$ points are distinct normalised projective points, and checks all
$\binom{20}{3}=1140$ triples by \emph{both} line membership and rank over
$\F_{3}$, obtaining $0$ collinear triples in each case; it also verifies
$20>14$, that the $20$-cap is inclusion-maximal, the $47$-point cap in
$\PG(4,5)$, the $16$-point cap in $\PG(4,2)$, and the quadric/formula mismatch.
It prints \texttt{PASS} and exits $0$ exactly when every check holds:

\begin{quote}\ttfamily
\$ python3 reproduce.py\\
... PASS: all checks verified; conjecture 00000001003 is FALSE.
\end{quote}

\section{Formalisation}

The refutation at $q=3$ is formalised in the Lean~4 project under
\texttt{lean4/} (core Lean only, \texttt{import Std}, no Mathlib, no
\texttt{sorry}, no \texttt{native\_decide}). Points are a structure
\texttt{Pt} with five \texttt{Nat} fields, arithmetic is reduced with
\texttt{\% 3}, and collinearity is decided by the executable predicate
\texttt{bad}, which is complete over $\F_{3}$: the projective line through two
distinct points $p,q$ is $\{p,q,\text{normalize}(p+q),\text{normalize}(p+2q)\}$.
The main statements are

\begin{quote}\ttfamily
def formula (q : Nat) : Nat := q*q + q + 1 + (q+1)/3\\
theorem pts20\_isCap : isCap pts20 = true\\
theorem formula\_three : formula 3 = 14\\
theorem refutes\_conjecture :\\
\hspace*{2em}$\exists$ S : List Pt, S.length = 20 $\wedge$ isCap S = true
  $\wedge$ formula 3 < S.length
\end{quote}

All eleven theorems are proved by \texttt{decide} and the audit
(\texttt{lake env lean Check.lean}, using \texttt{\#print axioms}) reports
``does not depend on any axioms'' for every one of them --- in particular no
\texttt{sorryAx} and no \texttt{Lean.ofReduceBool}. See
\texttt{lean4/README.md} for the statement table, proof strategy, and scope
note.

\section*{Status against the submission rules}

Rule~3 requires a LaTeX source, a PDF document, and a Lean~4 project.

\begin{itemize}
\item \textbf{LaTeX source} --- \texttt{main.tex}, a standalone \texttt{article}
  using \texttt{geometry}, \texttt{amsmath}, \texttt{amssymb}, \texttt{amsthm},
  \texttt{array}, \texttt{parskip}, \texttt{booktabs}; compiles with
  \texttt{tectonic --outdir build main.tex}.
\item \textbf{PDF} --- at \texttt{build/main.pdf}.
\item \textbf{Lean~4 project} --- \texttt{lean4/} with \texttt{lean-toolchain}
  pinned to \texttt{leanprover/lean4:v4.33.1}, \texttt{lakefile.toml} (library
  \texttt{Main}, name \texttt{tlmc1003}, no dependencies), \texttt{Main.lean},
  \texttt{Check.lean} (\texttt{\#print axioms}), and \texttt{lean4/README.md}.
  \texttt{lake build} exits $0$ and the audit reports no \texttt{sorryAx}.
\end{itemize}

\end{document}
